
\begin{abstract}
随着大语言模型在智能问答、内容生成和辅助决策等场景中的广泛应用，其输出中可能包含的群体刻板印象、差别化评价和不公平推断问题日益受到关注。对于闭源或仅提供推理接口的黑盒大语言模型而言，研究者通常无法访问模型参数和内部表示，这使得传统依赖模型内部信息的偏见分析与修正方法难以直接适用。与此同时，实际应用场景不仅需要识别模型是否产生了偏见输出，还需要在尽量不破坏原始语义和任务可用性的前提下，对已产生的偏见回复进行有效校正。针对上述问题，本文将研究黑盒大语言模型偏见修正方法，主要工作如下：

1. 针对黑盒大语言模型输出中存在公平性偏见的问题，提出一种基于反事实输入与思维链反思验证的偏见检测方法。该方法先对输入中的敏感属性进行抽取，并构造仅替换敏感属性而保持语义一致的反事实输入。然后，从原始回复与反事实回复之间的语义偏移、情感差异、毒性变化和生成偏好等多个维度构造候选偏见分数。最后，引导模型结合事实性与公平性，对原始回复与反事实回复之间的差异命题进行反思和验证。该方法无需访问模型参数和训练过程，能够从黑盒输入输出行为中识别由敏感属性变化触发的不合理回复差异。

2. 针对黑盒模型输出偏见修正中普遍存在的整体改写幅度大、语义损失严重的问题，提出一种基于优先级排序与最小编辑约束的偏见回复重写方法。首先对偏见回复进行句子级与短语级划分，并结合模型偏见评分、删除贡献分数和敏感属性关联分数定位需要处理的局部偏见片段。然后，引入编辑代价并构建重写优先级函数，对候选片段按照风险程度与修改成本进行排序。针对当前待处理片段生成局部候选替换片段，依据公平性分数、语义一致性分数和编辑代价综合选择最优重写结果。该方法将按照优先级顺序迭代执行上述局部重写过程。该方法能够在降低偏见风险的同时，尽量保持原始回复的语义、内容信息和任务可用性。

3. 在开放式文本与对话系统两类偏见评测数据集上开展实验，验证本文方法在偏见识别与偏见重写两项任务中的有效性。在偏见识别实验中，本文方法在多个评估指标与外部评审中均取得最优，平均偏见识别准确率达到 84.0\%，较三种现有识别方法分别提高 18.7\%、3.5\% 和 2.9\%。消融实验进一步表明，多维差异检测与反思验证机制分别使偏见识别准确率提高 7.8\% 和 1.5\%，两个模块均对整体性能提升具有积极作用。在偏见重写实验中，本文方法在语义保持度得分上更稳定地集中于高分区间，在校正效率上取得了更低的平均编辑距离比例和更高的平均风险降低程度，且偏见去除率比现有方法提高 2\%～3.8\%。实验结果表明，本文提出的方法能够较有效地识别黑盒大语言模型输出中的偏见，并在最小编辑约束下实现更优的偏见缓解效果。

\end{abstract}


\begin{abstractEn}
With the rapid deployment of large language models in intelligent question answering, content generation, and decision support, group stereotypes, differential evaluations, and unfair inferences in model outputs have become increasingly concerning. For black-box large language models that only provide inference interfaces, researchers usually cannot access model parameters or internal representations, which makes traditional bias analysis and mitigation methods that rely on internal model information difficult to apply directly. Meanwhile, practical applications require not only detecting whether a model output is biased, but also effectively revising biased responses while preserving the original semantics and task utility as much as possible. This thesis investigates a bias mitigation method for black-box large language models. The main contributions are as follows:

1. For bias detection in black-box large language model outputs, a method based on counterfactual inputs and chain-of-thought reflective verification is proposed. The method first extracts sensitive attributes from the input and constructs counterfactual inputs that only replace the sensitive attributes while preserving task semantics. It then builds candidate bias scores from multiple dimensions, including semantic shift, sentiment difference, toxicity change, and generation preference, between the original response and the counterfactual response. Finally, the model is guided to reflect on and verify the difference propositions between the original response and the counterfactual response by jointly considering factuality and fairness. This method does not require access to model parameters or training procedures, and can identify unreasonable response differences triggered by sensitive attribute changes from black-box input-output behavior.

2. To address the problems of large rewriting scope and severe semantic loss in bias mitigation for black-box model outputs, a bias response rewriting method based on priority ranking and minimum-edit constraints is proposed. The method first performs sentence-level and phrase-level segmentation on a biased response, and then locates local biased spans by combining model-based bias scores, deletion contribution scores, and sensitive-attribute association scores. Next, it introduces edit cost to construct a rewriting priority function that ranks candidate spans according to risk level and modification cost. Finally, for the current target span, local replacement candidates are generated, and the best rewritten response is selected by jointly considering fairness, semantic consistency, and edit cost, after which local rewriting is iteratively applied in priority order. This method reduces bias risk while preserving the semantics, content information, and task utility of the original response.

3. Experiments are conducted on two bias evaluation datasets, namely the open-ended text dataset and the dialogue-oriented dataset, to validate the effectiveness of the proposed method in both bias detection and bias rewriting tasks. In the bias detection experiments, the proposed method achieves the best results on multiple evaluation metrics and external reviews, with an average bias detection accuracy of 84.0\%, outperforming three existing detection methods by 18.7\%, 3.5\%, and 2.9\%, respectively. Ablation studies further show that multidimensional difference detection and reflective verification improve bias detection accuracy by 7.8\% and 1.5\%, respectively, indicating that both modules contribute positively to the overall performance. In the bias rewriting experiments, the proposed method exhibits more stable high-score distributions in semantic preservation, achieves lower average edit-distance ratios and higher average risk reduction in correction efficiency, and improves the bias removal rate by 2\% to 3.8\% over existing methods. The results demonstrate that the proposed method can effectively identify bias in black-box large language model outputs and achieve better bias mitigation under minimum-edit constraints.

\end{abstractEn}
